\section*{Bab \ref{ch:g6:shapes} --- Segitiga dan Segi Empat}

\begin{solution}{exo:g6:shapes:1}
Ikutilah caranya: \omterm{def:g6:lines:objects}{ruas garis} $[AB]$ sepanjang $6$ cm, busur
berjari-jari $5$ cm berpusat di $A$, busur berjari-jari $4$ cm
berpusat di $B$; titik potongnya adalah $C$.
\end{solution}

\begin{solution}{exo:g6:shapes:2}
Alas $[EF]$ sepanjang $3$ cm, lalu dua busur berjari-jari $5$ cm
berpusat di $E$ dan $F$; perpotongannya adalah $D$. Kedua sudut
alasnya $\widehat{DEF}$ dan $\widehat{DFE}$ sama besar (kira-kira
$72.5^\circ$ masing-masing) --- akibat simetri \omterm{def:g6:shapes:triangles}{segitiga sama kaki}
(\cref{ex:g6:symmetry:proves}).
\end{solution}

\begin{solution}{exo:g6:shapes:3}
\omterm{def:g6:lines:objects}{Ruas garis} sepanjang $4$ cm, lalu dua busur berjari-jari $4$ cm
berpusat di kedua ujungnya. Setiap sudutnya berukuran $60^\circ$.
\end{solution}

\begin{solution}{exo:g6:shapes:4}
(a) sama sisi (karena itu juga sama kaki). \quad
(b) sama kaki ($5 = 5$). \quad
(c) siku-siku dan sama kaki (dua sudut sama besar $45^\circ$
menghadap dua sisi yang sama panjang). \quad
(d) sama sisi (tiga sudut sama besar).
\end{solution}

\begin{solution}{exo:g6:shapes:5}
Sisi miring adalah sisi yang berhadapan dengan \omterm{def:g3:shapes:rightangle}{sudut siku-siku}, dan
selalu sisi yang terpanjang. Pada $MNP$ yang siku-siku di $N$, sisi
miringnya $[MP]$: ia tidak pernah menyentuh titik \omterm{def:g3:shapes:rightangle}{sudut siku-sikunya}, $N$.
\end{solution}

\begin{solution}{exo:g6:shapes:6}
\emph{1. Benar}: persegi punya empat \omterm{def:g3:shapes:rightangle}{sudut siku-siku}, dan itulah
definisi persegi panjang.

\emph{2. Salah}: persegi panjang $5 \times 3$ punya sisi yang tidak sama panjang.

\emph{3. Benar}: persegi punya empat sisi sama panjang, dan itulah
definisi \omterm{def:g6:shapes:quadrilaterals}{belah ketupat}.

\emph{4. Benar}: lihat \cref{ex:g6:shapes:recognize}.
\end{solution}

\begin{solution}{exo:g6:shapes:7}
Lukisannya memakai \omterm{def:g6:lines:perp}{garis tegak lurus}, seperti pada
\cref{ex:g6:shapes:program}. Kedua diagonalnya berukuran kira-kira
$5.8$ cm: diagonal persegi panjang sama panjang
(\cref{prop:g6:shapes:diagonals}).
\end{solution}

\begin{solution}{exo:g6:shapes:8}
Gambarlah \omterm{def:g6:lines:objects}{ruas garis} sepanjang $6$ cm; titik tengahnya $O$; \omterm{def:g6:lines:perp}{garis
tegak lurus} di $O$; padanya, dua titik berjarak $2$ cm dari $O$ di
kedua sisinya. Menghubungkan keempat ujungnya memberi \omterm{def:g6:shapes:quadrilaterals}{belah ketupat}
(keempat \omterm{def:g5:solids:def}{titik sudutnya} berjarak $3$ cm dan $2$ cm dari $O$ sepanjang
garis yang \omterm{def:g6:lines:perp}{tegak lurus}, jadi keempat sisinya sama panjang menurut simetri).
\end{solution}

\begin{solution}{exo:g6:shapes:9}
Satu program yang benar:
1. gambarlah \omterm{def:g6:lines:objects}{ruas garis} $[BC]$ dengan $BC = 4$ cm;
2. gambarlah busur berpusat $B$ dengan \omterm{def:g3:shapes:shapes}{jari-jari} $6$ cm di atas $[BC]$;
3. gambarlah busur berpusat $C$ dengan \omterm{def:g3:shapes:shapes}{jari-jari} $6$ cm di atas $[BC]$;
4. namailah $A$ titik potong kedua busur itu;
5. gambarlah \omterm{def:g6:lines:objects}{ruas garis} $[AB]$ dan $[AC]$.
\end{solution}

\begin{solution}{exo:g6:shapes:10}
Bagaimanapun letak $M$ pada lingkaran, sudut $\widehat{AMB}$ berukuran
$90^\circ$: segitiga $AMB$ siku-siku di $M$ setiap kali $[AB]$ menjadi
\omterm{def:g4:geometry:circle}{diameter}. Pembuktiannya menunggu \cref{ch:g8:cosine}.
\end{solution}

\begin{solution}{exo:g6:shapes:11}
Diagonal yang saling memotong di titik tengahnya dan sama panjang:
itulah daftar sifat \emph{persegi panjang}
(\cref{prop:g6:shapes:diagonals}). Tidak ada yang memaksa diagonalnya
\omterm{def:g6:lines:perp}{tegak lurus}, jadi ia tidak harus berupa \omterm{def:g6:shapes:quadrilaterals}{belah ketupat} atau persegi ---
persegi panjang adalah jawaban yang tepat (persegi juga memenuhinya,
tetapi tidak dipaksakan).
\end{solution}

\begin{solution}{exo:g6:shapes:12}
Dengan $4 \times 4$ titik ada $3 \times 3$ petak. Persegi yang \omterm{def:g6:lines:perp}{sejajar}
petak: sisi $1$, $2$ atau $3$ memberi $9 + 4 + 1 = 14$ persegi.
Persegi miring juga ada: yang kecil, yang sisinya diagonal petak
$1 \times 1$ (seperti $(1,0)$, $(2,1)$, $(1,2)$, $(0,1)$) --- ada
empat; dan yang lebih besar, yang sisinya diagonal persegi panjang
$2 \times 1$, seperti $(1,0)$, $(3,1)$, $(2,3)$, $(0,2)$ --- ada dua.
Keempat sisi persegi semacam itu sama panjang karena semuanya diagonal
persegi panjang yang identik. Seluruhnya: $14 + 4 + 2 = 20$ persegi.
(Hitungan mana pun yang menyertakan satu contoh miring yang benar dan
$14$ persegi \omterm{def:g6:lines:perp}{sejajar} petak sudah merupakan keberhasilan pada tingkat ini.)
\end{solution}

\begin{solution}{pb:g6:shapes:1}
\textbf{1.} Keempat sisinya berukuran sama (kira-kira
$4.2$ cm) dan keempat sudutnya siku-siku: bangunnya sebuah
\emph{persegi}. (Diagonal yang sama panjang memberi \omterm{def:g3:shapes:rightangle}{sudut siku-siku},
seperti pada persegi panjang; diagonal yang \omterm{def:g6:lines:perp}{tegak lurus} memberi sisi
yang sama panjang, seperti pada \omterm{def:g6:shapes:quadrilaterals}{belah ketupat}; keduanya sekaligus
memberi persegi, yaitu \cref{prop:g6:shapes:diagonals} yang dibaca terbalik.)

\textbf{2.} Sebuah \emph{\omterm{def:g6:shapes:quadrilaterals}{belah ketupat}}: empat sisi sama panjang
(kira-kira $4.7$ cm), tetapi tanpa \omterm{def:g3:shapes:rightangle}{sudut siku-siku} --- ini persis
lukisan pada \cref{exo:g6:shapes:8}.

\textbf{3.} Keempat sudutnya ternyata siku-siku tetapi sisinya
terdiri atas dua panjang yang berbeda: sebuah \emph{persegi panjang}
--- inilah keadaan pada \cref{exo:g6:shapes:11}.

\textbf{4.} Tidak ada pernyataan sama panjang yang bertahan setelah
diukur dan tidak ada \omterm{def:g3:shapes:rightangle}{sudut siku-siku} yang muncul; tetapi sisi yang
berhadapan sama panjang dua-dua dan tampak \omterm{def:g6:lines:perp}{sejajar}. Inilah
``persegi panjang yang miring'' pada umumnya, yang diagonalnya sekadar
saling membagi dua --- yaitu \emph{jajargenjang}, yang dinamai dan
dipelajari pada \cref{ch:g7:central}.

\textbf{5.} Diagonal yang berpotongan di titik tengah bersamanya:
\begin{center}
\renewcommand{\arraystretch}{1.2}
\begin{tabular}{l|l}
diagonalnya\dots & bangunnya \\
\hline
sama panjang dan \omterm{def:g6:lines:perp}{tegak lurus} & persegi \\
\omterm{def:g6:lines:perp}{tegak lurus} saja & \omterm{def:g6:shapes:quadrilaterals}{belah ketupat} \\
sama panjang saja & persegi panjang \\
bukan keduanya & si tanpa nama (lihat \cref{ch:g7:central}) \\
\end{tabular}
\end{center}

\textbf{6.} Setiap segitiga pojoknya punya satu sisi siku-siku pada
tiap diagonal: setengah dari $8$ cm dan setengah dari $5$ cm, jadi
sisi siku-sikunya $4$ cm dan $2.5$ cm, dengan \omterm{def:g3:shapes:rightangle}{sudut siku-siku} di
antaranya --- data yang sama untuk keempatnya. Segitiga yang dibangun
dari dua sisi siku-siku yang sama beserta \omterm{def:g3:shapes:rightangle}{sudut siku-siku} di antaranya
adalah salinan persis satu sama lain (bangunlah: bukaan jangkanya
identik). Karena itu keempat sisi miringnya --- yaitu sisi bangun itu
--- sama panjang: bangunnya \omterm{def:g6:shapes:quadrilaterals}{belah ketupat}.

\textbf{7.} Bukalah jangka sepanjang satu sisi, lalu jalankan ia
mengelilingi bangun itu: kalau jarum dan pensilnya mendarat tepat pada
tiap pasang pojok yang bertetangga, keempat sisinya sama dengan bukaan
itu. Tidak perlu penggaris.

\textbf{8.} Programnya:
\begin{enumerate}
  \item Gambarlah \omterm{def:g6:lines:objects}{ruas garis} $[AC]$ dengan $AC = 6$ cm lalu letakkan
        titik tengahnya $O$ (ukurlah $3$ cm dari $A$).
  \item Gambarlah \omterm{def:g6:lines:perp}{garis tegak lurus} terhadap $(AC)$ di $O$.
  \item Padanya, letakkan $B$ dan $D$ berjarak $3$ cm dari $O$, satu
        di tiap sisi.
  \item Hubungkan $A$, $B$, $C$, $D$.
\end{enumerate}
Diagonalnya sama panjang ($6$ cm), \omterm{def:g6:lines:perp}{tegak lurus}, dan saling membagi
dua: menurut kamusnya, $ABCD$ adalah persegi.

\textbf{9.} Pengukurannya memberi kira-kira $4.2$ cm: \emph{lebih
panjang} daripada $3$ cm --- terkaan yang lazim, ``setengah
diagonalnya'', ternyata salah. Setengah jawaban yang dapat dibuktikan:
untuk pergi dari pojok $A$ ke pojok $B$ berikutnya, kita dapat
berjalan lurus pada sisi $[AB]$, atau memutar lewat pusatnya
$A \to O \to B$ yang memakan $3 + 3 = 6$ cm; sisi yang lurus lebih
pendek, jadi sisinya berukuran kurang daripada $6$ cm. Bahwa ia juga
\emph{lebih} daripada $3$ cm adalah yang ditunjukkan penggaris;
\cref{ch:g8:pythagoras} akan membuktikannya dan memberi nilai
tepatnya ($3$ cm dikalikan akar kuadrat $2$, kira-kira $4.24$ cm).

\textbf{10.} Keempat segitiga pojok di $H$ sekarang terdiri atas
\emph{dua} pasang yang sepadan: dua dengan sisi siku-siku $3$ cm dan
$2.5$ cm (sisi miringnya $[AB]$ dan $[AD]$), dua dengan sisi siku-siku
$5$ cm dan $2.5$ cm (sisi miringnya $[CB]$ dan $[CD]$). Jadi $AB = AD$
dan $CB = CD$: \omterm{pb:g6:shapes:1}{layang-layang} punya dua pasang sisi \emph{bertetangga}
yang sama panjang (\omterm{def:g6:shapes:quadrilaterals}{belah ketupat} adalah \omterm{pb:g6:shapes:1}{layang-layang} istimewa yang
keempat sisinya sepadan).

\textbf{11.} \omterm{def:g4:geometry:polygon}{Segi empat}: $2$ diagonal. \omterm{def:g4:geometry:polygon}{Segi lima}: $5$.
\omterm{def:g4:geometry:polygon}{Segi enam}: $9$.

\textbf{12.} Dari satu \omterm{def:g5:solids:def}{titik sudut} \omterm{def:g4:geometry:polygon}{segi enam}, diagonalnya menuju
setiap \omterm{def:g5:solids:def}{titik sudut} kecuali dirinya sendiri dan kedua tetangganya:
$6 - 3 = 3$ diagonal. Enam \omterm{def:g5:solids:def}{titik sudut} meluncurkan
$6 \times 3 = 18$ ujung diagonal --- tetapi setiap diagonal sudah
terbilang pada \emph{kedua} ujungnya, sekali dari tiap ujungnya,
persis seperti jabat tangan yang dibilang oleh kedua orang yang
berjabat (\cref{pb:g6:wholes:1}). Hitungan yang sebenarnya:
$18 \div 2 = 9$. Itu cocok dengan gambarnya.

\textbf{13.} Sepuluh \omterm{def:g5:solids:def}{titik sudut}: $10 - 3 = 7$ dari masing-masing,
jadi $10 \times 7 \div 2 = 35$ diagonal. Dua puluh \omterm{def:g5:solids:def}{titik sudut}:
$20 \times 17 \div 2 = 170$ diagonal.

\textbf{14.} Cobalah beberapa calon: $10$ \omterm{def:g5:solids:def}{titik sudut} memberi $35$
(pertanyaan 13); $11$ \omterm{def:g5:solids:def}{titik sudut} memberi $11 \times 8 \div 2 = 44$.
Ketemu: \omterm{def:g4:geometry:polygon}{segi banyak} itu punya $11$ \omterm{def:g5:solids:def}{titik sudut}.

\textbf{15.} Segitiga sama sekali tidak punya diagonal: yang ``bukan
tetangga'' dari setiap \omterm{def:g5:solids:def}{titik sudutnya} adalah\,\dots\ tidak seorang
pun, karena kedua \omterm{def:g5:solids:def}{titik sudut} yang lain sama-sama tetangganya. Caranya
sepakat: $3 - 3 = 0$ dari setiap \omterm{def:g5:solids:def}{titik sudut}, seluruhnya $0$. Jabat
tangan dan diagonal adalah hitungan yang sama dengan pakaian yang
berbeda: pasangan yang dipilih dari sebuah kelompok, tiap pasangan
dibilang sekali --- matematika mendaur ulang gagasannya yang bagus.

\end{solution}
